\subsection{门控边界、污染交互与可追溯失败案例}
\label{sec:sensitivity-interpretation}

观测数量改变了信息量，也改变了模型可用性。三折时间交错划分要求每个训练折
至少有6份观测，因此计数条件为$n-\lceil n/3\rceil\geq6$。在$n=8$的96例中，
第一折训练集只有5份观测，程序直接停止CV，完成折数为0，运动可用标记全部为假，
最终全部输出静态模型。该运动无污染单元的均值误差为11.5264 m，恒速稳健
对照为4.1988 m。这里的静态选择来自实现门槛，不能解释成观测证据支持目标静止。

在$n\geq12$的288例中，三折均完成且运动模型全部可用。此时仍存在污染相关的
选择变化：真实运动案例在名义污染率$0,0.1,0.25,0.4$下，选择静态的次数依次为
$4/36,8/36,25/36,27/36$，对应平均位置误差为
3.6955、7.3217、8.8526、10.8942 m。CV比较历史观测上的标准化距离预测损失，
最终任务却评价参考时刻的位置；从距离误差到当前位置误差的放大受轨迹参数、
时间分布和几何条件影响，两类风险没有自动保持大小顺序的保证。

\begin{table}[H]
  \centering
  \small
  \caption{三个实际失败案例的配对位置误差。案例均保留于正式统计，单位为m。}
  \label{tab:sensitivity-failures}
  \begin{tabular}{llrrr}
    \toprule
    案例 & 最终选择 & 所选方法 & 恒速稳健 & 静态WLS\\
    \midrule
    A：运动，$n=16,\epsilon=0.1$ & 恒速 & 64.2834 & 64.2834 & 6.1442\\
    B：运动，$n=16,\epsilon=0.25$ & 静态 & 23.0179 & 3.1658 & 22.8234\\
    C：静态，$n=12,\epsilon=0.4$ & 恒速 & 41.5503 & 41.5503 & 9.3065\\
    \bottomrule
  \end{tabular}
\end{table}

案例A为本次最大误差，求解器仍报告成功、Jacobian秩为4、没有活动数值边界；
最小奇异值约0.16293，条件数约128.22。CV允许并选择恒速，但配对恒速基线也有
相同误差，说明失败不能只归因于模型选择。满秩与优化终止不能保证估计落在真解
附近，也不能代替理论局部误差界的邻域和噪声条件。案例B中运动模型可用，但CV
静态／运动得分为3.20774／3.51454，静态选择比配对恒速结果多出19.8521 m误差。
案例C则把零速度目标解释成运动，直接展示了静态非劣目标仍未得到保证。

为便于复核，表中A、B、C依次对应下列完整\texttt{case\_id}；在
\nolinkurl{results/sensitivity/predictions.json}中按该字段查找即可获得三种方法的
预测和CV诊断，在同目录\texttt{inputs.json}中可核对生成输入：
\begin{description}[leftmargin=1.6em,labelsep=0.4em,font=\normalfont\bfseries]
  \item[A] \nolinkurl{sensitivity-moving_outliers_n16_c1000-r05-921984553727470929}
  \item[B] \nolinkurl{sensitivity-moving_outliers_n16_c2500-r12-13744984418704934917}
  \item[C] \nolinkurl{sensitivity-static_outliers_n12_c4000-r04-14729045951357002381}
\end{description}

本分析为32格、每格12例的探索性描述。不同网格独立抽样，污染数量还经过
$\operatorname{round}(n\epsilon)$取整；观察到的交互和极端案例不能单独识别因果
机制，也不能精确定位污染崩溃阈值。上述解释结合代码规则与保存的诊断提出，
没有据此修改冻结方法、作逐格显著性宣称或建立新的覆盖保证。
